\section{Inverted-LiRA scores for destroyed-model FedRecovery variants}
\label{app:inverted-lira}

The LiRA attack with a shared shadow pool scores a target model's
per-sample loss against the shadow pool's in-set / out-set loss
distributions and reports the resulting likelihood ratio. For a model
with meaningful utility, the shadow pool's loss distributions have
positive separation between in-set and out-set samples, and the attack
reports an advantage based on how well the target model fits the in-set
distribution. For a model whose outputs have been destroyed by DP noise,
the per-sample losses are approximately uniform, which produces a
degenerate loss distribution. The shadow pool's distributions remain
calibrated to a non-destroyed model, so the destroyed-model losses sit
far from both the in-set and out-set means, in a direction
anti-correlated with the shadow pool's structure.

The result is an inverted AUROC ($<\!0.5$) and a corresponding
$|2 \cdot \mathrm{AUROC} - 1|$ advantage number that reflects the
distributional mismatch, not a successful membership-inference attack.

\begin{table}[h]
\centering
\small
\begin{tabular}{lrr}
\toprule
Method & LiRA AUROC & $|2\mathrm{AUROC}-1|$ (inverted) \\
\midrule
FedRecovery $\varepsilon = 1$     & $0.070$ & $0.859$ \\
FedRecovery $\varepsilon = 2$     & $0.393$ & $0.215$ \\
FedRecovery $\varepsilon = 5$     & $0.452$ & $0.106$ \\
FedRecovery $\varepsilon = 10$    & $0.484$ & $0.075$ \\
\bottomrule
\end{tabular}
\caption{Raw LiRA output for destroyed-model FedRecovery variants (mean across three seeds). AUROC $<0.5$ indicates an inverted score distribution; these values are attack artefacts, not comparable to Table~\ref{tab:cifar10-ci-baselines}.}
\label{tab:inverted-lira}
\end{table}
